\section{Proofs of the main results}
\label{sec:proofs}

The two theorems of this section share the same algebraic skeleton ---
both reduce a question about mutual information to a statement about
the joint law of two modular coordinates --- but they apply to two
different measures and have, accordingly, two different conclusions.
Theorem~\ref{thm:empirical_invariance} states that the mutual
information under the empirical measure is independent of the
parameter~$\mu$ that selects the Bianchi-I geometry $G(\mu)$;
Theorem~\ref{thm:ruelle_zero} states that the mutual information
under the Ruelle measure is exactly zero. Together they pin down the
arithmetic origin of the residual
$\mathcal{I}_{p,q}$: it is independent of any geometric reparametrisation
of the active primes (Theorem~\ref{thm:empirical_invariance}) and it
vanishes for the idealised Ruelle measure
(Theorem~\ref{thm:ruelle_zero}), so it must be attributable to the
arithmetic structure of the gap field $\{g_n\}$ alone. The closed-form
expression in Section~\ref{sec:closed_form} confirms this from a
constructive angle.

\subsection{Theorem~4.1: Geometric Decoupling Theorem (empirical form)}
\label{ssec:thm_4_1_proof}

\begin{theorem}[Geometric Decoupling Theorem, empirical form]
\label{thm:empirical_invariance}
\footnote{Formally verified in Lean~4 + Mathlib as
\texttt{PT.CrtDecoupling.Empirical.empirical\_invariance}
(hypothesis form),
\texttt{SpectralReduction.empirical\_invariance\_closed}
(finite-state closed form),
\texttt{Phase3.empirical\_invariance\_infinite}
(infinite-state form via Birkhoff--Hopf), and
\texttt{Phase4.sieve\_empirical\_invariance}
(sieve-specific closure); see Section~\ref{ssec:formal_verification}.}
Let $p, q \in \{3, 5, 7\}$ be distinct active primes, and let
$\Phi_p \in \mathcal{A}_p$, $\Phi_q \in \mathcal{A}_q$ be bounded
observables on the corresponding modular $\sigma$-algebras. The
function
\begin{equation}
  \mathcal{I}_{p,q}(\mu)
  \;:=\;
  I_{\pi_m^{\mathrm{emp}}}\!\bigl(\Phi_p \,;\, \Phi_q \,\big|\, G(\mu)\bigr)
  \label{eq:Ipq_def}
\end{equation}
is constant on $\mu \in [\mu_c, \infty)$.
\end{theorem}

The conditional mutual information in~\eqref{eq:Ipq_def} is taken in
the sense of Cover and Thomas~\cite[\S~2.5]{CoverThomas2006}: writing
$G$ for a random variable taking the value $G(\mu)$, we have
$I(\Phi_p; \Phi_q \mid G) = \mathbb{E}\bigl[I_{G}(\Phi_p; \Phi_q)\bigr]$,
where $I_{G}$ denotes the mutual information under the conditional
law given $G$. For our purposes, the only fact we need is the standard
identity~\cite[Eq.~(2.66)]{CoverThomas2006}: if $G$ is $\mu$-a.s.\
constant under the measure of interest, then
$I(\Phi_p; \Phi_q \mid G) = I(\Phi_p; \Phi_q)$.

\begin{proof}
The proof proceeds in four steps.

\textbf{Step 1 (Empirical measure is geometry-independent).}
The gap sequence $\{g_n\}$ is defined by the Eratosthenes sieve, a
purely arithmetic algorithm that takes no reference to any continuous
parameter $\mu$ nor to any metric on a continuum. Consequently the
empirical stationary measure
$\pi_m^{\mathrm{emp}}$ in Definition~\ref{def:pi_emp} is a functional
of $\{g_n\}$ alone, and is independent of both~$\mu$ and the family
$G(\mu)$. This is the content of Remark~\ref{rmk:residual_arithmetic}
in the setup.

\textbf{Step 2 ($G(\mu)$ is spectrally measurable).}
By Lemma~\ref{lem:G_measurable} of Section~\ref{sec:geometric_fiber},
each component of $G(\mu)$ is a measurable functional of the spectrum
of $T_m$:
\[
  G(\mu) \;\in\; \sigma\bigl(\mathrm{Spec}(T_m)\bigr)
  \qquad \text{for every } \mu \in [\mu_c, \infty).
\]
Since $\mathrm{Spec}(T_m)$ is fixed by the modular reduction
$r_n^{(m)} \mapsto r_{n+1}^{(m)}$ --- it is a property of the matrix
$T_m$ on $\mathcal{S}_m$, not of the trajectory ---
$\sigma(\mathrm{Spec}(T_m))$ is contained in the shift-invariant
$\sigma$-algebra $\mathcal{I} = \{A \in \mathcal{F} : T^{-1}A = A\}$
on the path space $\Omega$. We therefore have
$G(\mu) \in \mathcal{I}$.

\textbf{Step 3 (Ergodicity $\Rightarrow$ a.s.\ constancy).}
Proposition~\ref{prop:ergodic} of Section~\ref{sec:ergodic} states
that the shift $T$ on $\Omega$ is ergodic with respect to
$\pi_m^{\mathrm{emp}}$, by the Birkhoff--Hopf ergodic
theorem~\cite[Thm.~1.6]{Walters1982} applied to the spectral-gap
estimate of Theorem~T4 in~\cite[Ch.~7]{PT_Monograph}.
Corollary~\ref{cor:trivial} then asserts that the invariant
$\sigma$-algebra is $\pi_m^{\mathrm{emp}}$-trivial:
$\mathcal{I} = \{\emptyset, \Omega\} \pmod{\pi_m^{\mathrm{emp}}}$.
Combined with Step~2, this gives
\[
  G(\mu) \in \mathcal{I}
  \;\Longrightarrow\;
  G(\mu) \text{ is } \pi_m^{\mathrm{emp}}\text{-a.s.\ constant},
\]
which is Corollary~\ref{cor:G_constant}.

\textbf{Step 4 (Conditioning on a constant is trivial).}
By the standard identity recalled above the proof
(cf.~\cite[\S~2.5]{CoverThomas2006}), conditioning the joint law of
$(\Phi_p, \Phi_q)$ on a random variable that is almost surely
constant under the relevant measure does not change the joint
information content:
\begin{equation}
  \mathcal{I}_{p,q}(\mu)
  \;=\; I_{\pi_m^{\mathrm{emp}}}(\Phi_p\,;\,\Phi_q \mid G(\mu))
  \;=\; I_{\pi_m^{\mathrm{emp}}}(\Phi_p\,;\,\Phi_q).
  \label{eq:Ipq_reduced}
\end{equation}
The right-hand side depends only on the gap field and the modular
$\sigma$-algebras $\mathcal{A}_p, \mathcal{A}_q$. Since both of these
objects are independent of~$\mu$ (Step~1),
$\mathcal{I}_{p,q}(\mu)$ is independent of~$\mu$.
\end{proof}

\paragraph{What the proof is, and what it is not.}
The mathematical content of Steps~1--4 is fairly mild: an
arithmetically-defined object cannot depend on a continuous parameter
that does not enter its construction, and conditioning on a constant
removes the conditioning. What the proof does \emph{not} say is that
the inter-channel mutual information~$\mathcal{I}_{p,q}$ vanishes; it
only says that, whatever value $\mathcal{I}_{p,q}$ takes, it is the
same value at every $\mu \in [\mu_c, \infty)$. The actual numerical
value is computed in Section~\ref{sec:closed_form}
(Corollary~\ref{cor:closed_form}), where the closed form is fitted to
the corpus values $0.138$ bits and $0.283$ bits of
Remark~\ref{rmk:marginal_vs_joint}. The non-trivial epistemic claim
is therefore not that $\mathcal{I}_{p,q}$ is computable, but that it
is \emph{geometric-invariant} despite the appearance of $G(\mu)$ in
the conditioning of~\eqref{eq:Ipq_def}: the emergent geometry does
not couple the modular fibres. We return to this point in
Section~\ref{sec:applications}.

\subsection{Theorem~4.2: Geometric Decoupling Theorem (Ruelle form)}
\label{ssec:thm_4_2_proof}

\begin{theorem}[Geometric Decoupling Theorem, Ruelle form]
\label{thm:ruelle_zero}
\footnote{Formally verified in Lean~4 + Mathlib as
\texttt{PT.CrtDecoupling.Main.geometric\_decoupling\_ruelle};
see Section~\ref{ssec:formal_verification}.}
Under the Ruelle measure $\pi_m^{\mathrm{Ruelle}}$ of the
transfer matrix $T_m$, for all distinct active primes $p, q$ and
bounded observables $\Phi_p, \Phi_q$:
\begin{equation}
  I_{\pi_m^{\mathrm{Ruelle}}}\!\bigl(\Phi_p \,;\, \Phi_q\bigr)
  \;=\; 0
  \qquad (\text{exact}).
  \label{eq:ruelle_zero}
\end{equation}
\end{theorem}

\begin{proof}
By Theorem~\ref{thm:crt_tensor} of
Section~\ref{sec:crt_factor}, the CRT identification
$\mathbb{Z}/pq\mathbb{Z} \cong \mathbb{Z}/p\mathbb{Z} \times
\mathbb{Z}/q\mathbb{Z}$ lifts to the tensor decomposition
$T_{pq} = T_p \otimes T_q$. By Lemma~\ref{lem:pi_factor} of the same
section, the unique left Perron eigenvector
$\pi_{pq}^{\mathrm{Ruelle}}$ of $T_{pq}$ for the simple eigenvalue
$\lambda_0 = 1$ factorises as
\begin{equation}
  \pi_{pq}^{\mathrm{Ruelle}}
  \;=\; \pi_{p}^{\mathrm{Ruelle}} \otimes \pi_{q}^{\mathrm{Ruelle}},
  \label{eq:pi_factor_recall}
\end{equation}
where uniqueness of the Perron eigenvalue (and hence of the
eigenvector) on each factor is guaranteed by the spectral gap of
$T_p$ established in~\eqref{eq:Tp_spectrum}.
Equation~\eqref{eq:pi_factor_recall} is exactly the independence
condition for the joint law of $(r^{(p)}, r^{(q)})$ on
$\mathcal{S}_p \times \mathcal{S}_q$, so for any bounded
observables $\Phi_p \in \mathcal{A}_p$, $\Phi_q \in \mathcal{A}_q$,
\begin{equation}
  D_{\mathrm{KL}}\!\bigl(\pi_{pq}^{\mathrm{Ruelle}}
       \,\|\, \pi_{p}^{\mathrm{Ruelle}} \otimes \pi_{q}^{\mathrm{Ruelle}}\bigr)
  \;=\; 0,
  \label{eq:KL_zero}
\end{equation}
and the mutual information
$I_{\pi_{pq}^{\mathrm{Ruelle}}}(\Phi_p; \Phi_q) =
D_{\mathrm{KL}}\bigl(\mathbb{P}_{\Phi_p, \Phi_q} \,\|\,
\mathbb{P}_{\Phi_p} \otimes \mathbb{P}_{\Phi_q}\bigr)$ inherits the
vanishing from~\eqref{eq:KL_zero} by the data-processing
inequality~\cite[\S~2.8]{CoverThomas2006}: passing from
$(r^{(p)}, r^{(q)})$ to the deterministic image
$(\Phi_p, \Phi_q)$ cannot increase information content, and the
source already has zero mutual information.
\end{proof}

\paragraph{Numerical confirmation of~\eqref{eq:pi_factor_recall}.}
The factorisation~\eqref{eq:pi_factor_recall} is verified numerically
in Appendix~\ref{app:numerics}. With $m = pq = 15$, the script
\texttt{scripts/test\_crt\_decoupling.py} computes the left Perron
eigenvector $\pi_{15}^{\mathrm{Ruelle}}$ of $T_{15}$ directly by
solving $\pi T_{15} = \pi$ at \texttt{mpmath} precision
$50$~digits, and compares it to the Kronecker product
$\pi_3^{\mathrm{Ruelle}} \otimes \pi_5^{\mathrm{Ruelle}}$. The
observed discrepancy is
\begin{equation}
  \bigl\| \pi_{15}^{\mathrm{Ruelle}}
        - \pi_3^{\mathrm{Ruelle}} \otimes \pi_5^{\mathrm{Ruelle}}
  \bigr\|_2
  \;=\; 5.88 \times 10^{-53},
  \label{eq:T2_value}
\end{equation}
i.e.\ machine-precision zero at $50$~dps. Sampling
$N = 10^6$ i.i.d.~draws from
$\pi_3^{\mathrm{Ruelle}} \otimes \pi_5^{\mathrm{Ruelle}}$ and
estimating the empirical mutual information yields
$1.45 \times 10^{-6}$ bits, consistent with the
$O(1/\sqrt{N})$ statistical noise floor and not with a non-zero
signal. The two checks confirm Theorem~\ref{thm:ruelle_zero} both
algebraically (Eq.~\eqref{eq:T2_value}) and statistically.

\subsection{Remarks on the two theorems together}
\label{ssec:both_thms}

We collect three observations that orient the rest of the paper.

\paragraph{(i) Identifying the residual.}
Define the empirical residual on disjoint active channels by
\begin{equation}
  \mathcal{I}_{p,q}
  \;:=\;
  I_{\pi_m^{\mathrm{emp}}}\bigl(r^{(p)}\,;\, r^{(q)}\bigr)
  \;-\;
  I_{\pi_m^{\mathrm{Ruelle}}}\bigl(r^{(p)}\,;\, r^{(q)}\bigr).
  \label{eq:residual_def}
\end{equation}
Theorem~\ref{thm:ruelle_zero} forces the second term to vanish
identically, so $\mathcal{I}_{p,q} =
I_{\pi_m^{\mathrm{emp}}}(r^{(p)}; r^{(q)})$.
Theorem~\ref{thm:empirical_invariance} forces the first term to be
independent of the geometric parameter~$\mu$, so $\mathcal{I}_{p,q}$
is a single number per pair $(p, q)$. The corpus
values~\eqref{eq:MI_corpus} are
\begin{equation}
  \mathcal{I}_{3,5} \;=\; 0.138 \text{ bits},
  \qquad
  \mathcal{I}_{3,7} \;=\; 0.283 \text{ bits}.
  \label{eq:residual_values}
\end{equation}

\paragraph{(ii) The residual is arithmetic, not geometric.}
Step~1 of the proof of Theorem~\ref{thm:empirical_invariance}
identifies the empirical measure as a functional of the gap field
alone. The non-vanishing of $\mathcal{I}_{p,q}$ relative to the
Ruelle baseline must therefore be sourced by the arithmetic
constraints on consecutive prime gaps that the Ruelle Markov chain
does not see --- in particular by the multi-modular correlations
between $g_n \bmod p$ and $g_n \bmod q$ that survive once one
restricts attention to the trajectory $\omega^\ast$ generated by the
actual primes. The closed-form expression of
Corollary~\ref{cor:closed_form} in the next section provides a
quantitative model of this residual, fitted to the corpus
values~\eqref{eq:residual_values}.

\paragraph{(iii) Consequence for non-locality.}
The combination of the two theorems is the precise statement, at the
level of measure theory and information theory, that the emergent
Bianchi-I geometry $G(\mu)$ does not couple the disjoint modular
fibres of the prime sieve. Whatever inter-channel correlations the
empirical measure carries, they are arithmetic in origin and
unchanged by every choice of $\mu \in [\mu_c, \infty)$, including the
canonical fixed point $\mu^\ast = 15$. This is the technical input
required by the companion foundational
paper~\cite{PT_HOLONOMY_NONLOCALITY}, where it is used to argue that
the observed Bell-type correlations between space-like separated
detectors are inherited from the modular structure of the underlying
sieve, not transported by a spatial light-cone of the emergent
geometry. We say more about this in Section~\ref{sec:applications};
the technical lifting from the present finite-dimensional setting to
the infinite-dimensional CRT product $\mathcal{H}_\infty$ is the
content of Lemma~G in the
website corpus~\cite[Lemma G]{PT_Monograph}.
